% !PRINTED 316
\weberpage{316}

\begin{center}
\S.\ 104.

\gesperrt{Bestimmung einer oberen Grenze für die Wurzeln.}
\end{center}

Wir haben schon im dritten Abschnitt von dem Satze Gebrauch gemacht, dass
man immer eine positive Zahl finden kann, die grösser ist als die
absoluten Werthe sämmtlicher Wurzeln einer gegebenen Gleichung. Es kommt
aber jetzt darauf an, eine möglichst einfache und zugleich möglichst
genaue Bestimmung einer solchen oberen Grenze zu geben.

Wir betrachten zunächst nur die reellen Wurzeln, und suchen also eine
positive Zahl, die grösser ist als die grösste positive Wurzel einer
gegebenen Gleichung $f(x)=0$.

Eine solche Bestimmung giebt uns das \gesperrt{Budan}'sche Theorem.

Nehmen wir den Co\"efficienten der höchsten Potenz von $x$ in $f(x)$
gleich 1 (oder wenigstens positiv) an, so kann man die positive Zahl
$\alpha$ immer so gross annehmen, dass
\weq{1}{f(\alpha),\ f'(\alpha),\ f''(\alpha)\ldots f^{(n)}(\alpha)}
alle positiv werden. Dann geht in der Reihe der Functionen
\[
f(x),\ f'(x),\ f''(x)\ldots f^{(n)}(x)
\]
zwischen $x=\alpha$ und $x=\infty$ kein Zeichenwechsel mehr verloren, und
es kann also auch nach dem Theorem I keine Wurzel von $f(x)$ zwischen
$\alpha$ und $\infty$ liegen.

Will man aus dem gleichen Satze für die negativen Wurzeln eine untere
Grenze haben, so nehme man die positive Zahl $\beta$ so an, dass
\weq{2}{(-1)^n f(-\beta),\ (-1)^{n-1} f'(-\beta),\ \ldots f^{(n)}(-\beta)}
alle positiv werden, dann hat die Gleichung $f(x)=0$ sicher keine
negative Wurzel unter $-\beta$.

Diese Bestimmung der Grenzen rührt schon von \gesperrt{Newton} her.

Eine andere Bestimmung einer oberen Grenze, die für die Rechnung
einfacher ist, hat \gesperrt{Laguerre} angegeben\footnote{\gesperrt{Laguerre},
Nouvelles Annales d.\ math., 2.\ ser., t.\ XIX, 1880. Journal de math.,
3.\ ser., t.\ IX, 1883.}.

Er benutzt statt der Ableitungen $f(x), f'(x), f''(x)\ldots$ die
Functionen $f_0(x), f_1(x), f_2(x), \ldots f_{n-1}(x)$, die uns schon mehr-
%CONT>>
